\section{Conclusion}
\label{sec:conclusion}

In this paper, we introduced \texttt{LLMOrch-VAPT}, a resilient, multi-agent orchestration framework designed to address the operational and economic barriers to autonomous penetration testing. By decoupling security reasoning from specific LLM providers and implementing a provider-agnostic abstraction layer, we have demonstrated that autonomous red teaming engagements can be made robust against the volatility of the AI service landscape.

Our evaluation of 1,500 security reasoning tasks confirms that the three primary technical innovations introduced in this work—Autonomous Provider Failover (APF), Dynamic Complexity-Aware Tiering (DCAT), and Security-Semantic Caching (SSC)—collectively enable a high-performance, cost-effective VAPT infrastructure. Specifically, our system achieves an MTTR of less than 2 seconds during provider outages and maintains a 97.4\% reasoning success rate while reducing operational costs by 84.2\%. Furthermore, our 100\% adherence to human-in-the-loop safety gates ensures that these autonomous capabilities can be deployed responsibly within enterprise environments.

As LLMs continue to evolve, the need for intelligent, resilient orchestration will only increase. \texttt{LLMOrch-VAPT} provides the foundational infrastructure required to transition autonomous security agents from academic prototypes to production-grade defensive assets. By open-sourcing our implementation, we hope to foster a new generation of resilient, transparent, and cost-efficient autonomous security systems that can keep pace with the rapidly changing threat landscape.
